\subsection*{Training on more distinct translations}
To test whether the numeral transfer extended beyond familiar sentence frames, we varied the number of reviewed source passages while holding the number of training presentations fixed (Fig.~\ref{fig:translation-coverage}). A narrow writer trained on 108 distinct reviewed translations and a broad writer trained on 841, drawn from as many chapters. Both replayed all 55,872 original book pairs. On the 108 shared supervised passages, the narrow writer reproduced 49 reviewed answers exactly and the broad writer reproduced two. On the 733 additional passages, the broad writer reproduced seventeen, the narrow writer four and the starting writer five. One acquired answer modernises the biography of L\"u Yi while preserving his courtesy name and birthplace.

On 192 passages from unseen chapters, the broader training reduced edits against the archive gloss from 9,228 to 8,383. Literal copying incurred 5,612 edits. Neither trained writer conveyed a complete meaning on the eighteen designated readings. Each preserved the specified quantity on both sides of four of eight changes, while literal copying preserved all eight. Both learned writers attached the Chinese unit for years of age to a count representing the size of the celestial circuit. The local text teacher gave twelve faithful translations on eighteen readings and conveyed both complete meanings on five pairs. One method-masked reader assessed the answers. Its acceptance of nine of eighteen archive glosses reflects the difficulty of judging an isolated source fragment against a translation that sometimes supplies missing context. Broader supervision improved a surface measure in this comparison. It did not improve either registered semantic outcome.
